\documentclass[a4paper,11pt]{report}
\usepackage{amsmath}
\usepackage{amsfonts}
\usepackage{amssymb}
\usepackage{latexsym}
\usepackage{graphicx}
\usepackage{lineno}
\usepackage[T1]{fontenc}
\usepackage[utf8]{inputenc} % acentua??o
%\usepackage[brazil]{babel}
%\usepackage{amsfonts}
\usepackage{graphicx,graphpap}
\usepackage{enumitem}
%\graphicspath{{dissertação/}}
%graphpap permite quadricular figuras com o comando graphpaper para sabermos onde
%inserir as letras. depois de quadricularmos porcentamos o comando
\usepackage{blindtext}
\usepackage{multicol}
\usepackage{hyperref}
\usepackage{makeidx}
%\usepackage{eufrak}
%\usepackage{pst-plot}
\usepackage{graphicx,color}
\usepackage{float}
\usepackage[top=3cm,left=3cm,right=2cm,bottom=2cm]{geometry}
\usepackage{indentfirst}
\usepackage{epstopdf}
%\usepackage{yhmath}
\usepackage{fancyhdr}
\usepackage{anonchap}
\usepackage{verbatim}
\usepackage{algorithm}
\usepackage{algorithmic}
\usepackage{multirow}
\usepackage{caption}
\captionsetup[table]{skip=10pt}
\usepackage{subcaption}
\usepackage{xlop}
\usepackage{parcolumns}

\newtheorem{teorema}{Theorem}[chapter]
\newtheorem{corolario}[teorema]{Corollary}
\newtheorem{definicao}[teorema]{Definition}
\newtheorem{lema}[teorema]{Lemma}
\newtheorem{prop}[teorema]{Proposition}
\newtheorem{obs}[teorema]{Remark}
\newtheorem{exe}[teorema]{Example}
\newcommand{\whitebox}{$\Box$}
\newcommand{\blackbox}{$\blacksquare$}
\newtheorem{ex}[teorema]{Exerc\'icio}
\newtheorem{assump}{Assumption}
\renewcommand\theassump{A\arabic{assump}}
\newtheorem{scheme}{Scheme}

\newenvironment{prova}[1][Proof]{\textbf{#1.} }{\hfill\whitebox\bigskip}
%\newtheorem{obs}[Remark]{\textbf{#1.} }{\hfill\whitebox\bigskip}

\newcommand{\fracd}[2]{  \frac {\displaystyle {#1}}{\displaystyle {#2}}   }
%\newcommand{\intd}[3]{ \displaystyle {\int_{#1}^{#2}{#3}}d\theta}
\newcommand{\inte}[2]{ \displaystyle {\int_{#1}{#2}}dx}
\newcommand{\g }{\gamma}
\newcommand{\la }{\lambda}
\newcommand{\ep}{\epsilon}
\newcommand{\implica}{\Rightarrow}
\newcommand{\seta}{\rightarrow}
\newcommand{\esp}{\vspace*{0.5 cm}}
\newcommand{\h}{\hspace*{8mm}}
\newcommand{\divi}{\textrm{div}}
\newcommand{\R}{\mathbb{R}}
\newcommand{\demo}{{\bf{Demonstra\c{c}\~{a}o:}}\\}
\newcommand{\sen}{\textrm{sen}}
\newcommand{\senh}{\textrm{senh}}
\newcommand{\cupdot}{\mathbin{\mathaccent\cdot\cup}}
\newcommand{\Tau}{\mathcal{T}}
\newcommand{\bTau}{\partial\mathcal{T}}
\newcommand{\Xtil}{\widetilde{X}}
\newcommand{\Mtil}{\widetilde{M}}
\newcommand{\Xbar}{\bar{X}}
\newcommand{\V}{\mathcal{V}}
\newcommand{\Om}{\Omega}
\newcommand{\bOm}{\partial\Omega}
\newcommand{\bK}{\partial K}
\newcommand{\dt}{\partial_t}
\newcommand{\essup}{\textrm{ess\ sup}}
\newcommand{\Epsilon}{\mathcal{E}}
\newcommand{\n}{\mathbf{n}}
\newcommand{\ttau}{\tilde{\tau}}
\newcommand{\tlambda}{\tilde{\lambda}}
\newcommand{\haat}{\widehat}
\newcommand{\pr}{\mathcal{P}}
\newcommand{\prt}{\mathcal{P}^{\Delta t}}
\newcommand{\epsilont}{\mathcal{E}^{\Delta t}}
\newcommand{\etal}{\eta_{\lambda}}
\newcommand{\etali}{\eta_i}
\newcommand{\etatau}{\eta_{\tau}}
\newcommand{\hetatau}{\widehat{\eta_{\tau}}}
\newcommand{\etataui}{\eta_{\tau,i}}
\newcommand{\etaf}{\eta_f}
\newcommand{\hetaf}{\widehat{\eta_f}}
\newcommand{\etalik}{\eta_{\lambda,K,i}}
\newcommand{\etatauik}{\eta_{\tau,K,i}}
\newcommand{\Deltat}{\Delta t}
\newcommand{\DeltaT}{\Delta T}
\newcommand{\Pp}{\mathbb{P}}
\newcommand{\llambda}{\lambda^*}
\newcommand{\tiu}{\tilde{u}}
\newcommand{\tpsi}{\tilde{\psi}}
\newcommand{\Lam}{\Lambda^{\Deltat}_h}
\newcommand{\tLam}{\tilde{\Lambda}^{\Deltat}_h}
\newcommand{\muht}{\mu^{\Deltat}_h}
\newcommand{\tmuht}{\tilde{\mu}^{\Deltat}_h}
\newcommand{\hpsi}{\widehat{\psi}}
\newcommand{\hetali}{\widehat{\eta_{i}}}
\newcommand{\hetal}{\widehat{\eta_{\lambda}}}
\newcommand{\ddt}{\dfrac{d}{dt}}
\newcommand\norm[1]{\left\lVert#1\right\rVert}


\def\dash---{\kern.16667em---\penalty\exhyphenpenalty\hskip.16667em\relax}

%--------------cabeçalho do documento----------------
\fancyhf{}
\fancyhead[L]{\rightmark}
\fancyfoot[C]{\thepage}
\pagestyle{fancy}
%--------------------------------------------------------
\begin{document}
\begin{tabular}{lllll}
\multirow{3}*{\includegraphics[width=60px]{CFUVR/logo.png}}&\multicolumn{4}{c}{\large \textbf{Centro Federal de Educação Tecnológica de Minas Gerais}}\\
&\textbf{Professor}: Lucas Martins Rocha\ \  \textbf{Disciplina}: EDO \\
&\textbf{Turmas}: 1\ \ \textbf{Semestre}: 25-2 \\
&\textbf{Atividade}: Avaliação 2 \ \ \ \textbf{Nota}:\underline{\hspace{1cm}}/19pts\\
&\textbf{Estudante}:\underline{\hspace{10cm}}
\end{tabular}
%-------------------------------------------------------
\begin{itemize}
    \item Respostas sem justificativas \textbf{NÃO} serão consideradas.
    \item Essa prova tem o valor de \textbf{19 pontos}.

\end{itemize}

\begin{enumerate}
    \item Considerando a equação diferencial
    \[
    (x+3)y'' + (x+2)y' - y = 0,
    \]
    faça o que se pede:
    \begin{enumerate}
        \item Mostre que $y_1(x)=e^{-x}$ é ima solução da EDO. (2 pts)
        \item Determine uma segunda solução $y_2(x)$ de forma que $y_1$ e $y_2$ sejam soluções fundamentais da equação. (3 pts)
        \item Determine a solução geral da EDO e, em seguida, a solução do PVI com condição inciial $y(1)=1$ e $y'(1)=3$. (2 pts)
    \end{enumerate}
    \newpage
    \textbf{Estudante}:\underline{\hspace{10cm}}

    \item  (6 pts) Use o método de \textbf{variação dos parâmetros} para resolver o PVI
    \[
    \begin{cases}
        y'' + y' - 2y = t^2 + 3\\
        y(0)=0, y'(0)=0
    \end{cases}
    \]
    Resolva de forma organizada, deixando claro todo o processo de resolução. Não faça apenas cálculos, um texto orientando a leitura da resolução é essencial.
    
    \newpage
    \textbf{Estudante}:\underline{\hspace{10cm}}
    \item (6 pts) Use o método de \textbf{coeficientes a determinar} para resolver o PVI
    \[
    \begin{cases}
        y'' + 2y' + 2y = e^t\sen(t)\\
        y(0)=0, y'(0)=0
    \end{cases}
    \]
    Resolva de forma organizada, deixando claro todo o processo de resolução. Não faça apenas cálculos, um texto orientando a leitura da resolução é essencial.
    
\end{enumerate}

\end{document}